pdfoutput=1
% Options for packages loaded elsewhere
\PassOptionsToPackage{unicode}{hyperref}
\PassOptionsToPackage{hyphens}{url}
\pdfoutput=1
\documentclass[
  12pt,
]{article}
\usepackage{xcolor}
\usepackage{amsmath,amssymb}
\setcounter{secnumdepth}{-\maxdimen} % remove section numbering
\usepackage{iftex}
\ifPDFTeX
  \usepackage[T1]{fontenc}
  \usepackage{textcomp} % provide euro and other symbols
\else % if luatex or xetex
  \usepackage{unicode-math} % this also loads fontspec
  \defaultfontfeatures{Scale=MatchLowercase}
  \defaultfontfeatures[\rmfamily]{Ligatures=TeX,Scale=1}
\fi
\usepackage{lmodern}
\ifPDFTeX\else
  % xetex/luatex font selection
\fi
% Use upquote if available, for straight quotes in verbatim environments
\IfFileExists{upquote.sty}{\usepackage{upquote}}{}
\IfFileExists{microtype.sty}{% use microtype if available
  \usepackage[]{microtype}
  \UseMicrotypeSet[protrusion]{basicmath} % disable protrusion for tt fonts
}{}
\makeatletter
\@ifundefined{KOMAClassName}{% if non-KOMA class
  \IfFileExists{parskip.sty}{%
    \usepackage{parskip}
  }{% else
    \setlength{\parindent}{0pt}
    \setlength{\parskip}{6pt plus 2pt minus 1pt}}
}{% if KOMA class
  \KOMAoptions{parskip=half}}
\makeatother
\usepackage{longtable,booktabs,array}
\usepackage{calc} % for calculating minipage widths
% Correct order of tables after \paragraph or \subparagraph
\usepackage{etoolbox}
\makeatletter
\patchcmd\longtable{\par}{\if@noskipsec\mbox{}\fi\par}{}{}
\makeatother
% Allow footnotes in longtable head/foot
\IfFileExists{footnotehyper.sty}{\usepackage{footnotehyper}}{\usepackage{footnote}}
\makesavenoteenv{longtable}
\setlength{\emergencystretch}{3em} % prevent overfull lines
\providecommand{\tightlist}{%
  \setlength{\itemsep}{0pt}\setlength{\parskip}{0pt}}
\usepackage{bookmark}
\IfFileExists{xurl.sty}{\usepackage{xurl}}{} % add URL line breaks if available
\urlstyle{same}
\hypersetup{
  hidelinks,
  pdfcreator={LaTeX via pandoc}}

\author{SCX}
\date{2026}

\begin{document}

\section{Bayesian Update Interpretation of SCX
Self-Evolution}\label{bayesian-update-interpretation-of-scx-self-evolution}

\begin{quote}
\textbf{Version}: 2026-06-28 \textbar{} \textbf{Status}: Theoretical
framework \textbar{} \textbf{Scope}: Formal Bayesian analysis of the
self-evolution loop \textbf{Notation note}: In this document, \(S\)
denotes the state space (per existing SCX theory). The gatekeeper
scoring function at time \(t\) is denoted \(S_t\), distinguished by the
subscript \(t\). This follows the KEY NOTATION DISAMBIGUATION rule:
\(S_t\) = gatekeeper scoring function, \(S\) = state space.
\end{quote}

\begin{center}\rule{0.5\linewidth}{0.5pt}\end{center}

\subsection{Table of Contents}\label{table-of-contents}

\begin{enumerate}
\def\labelenumi{\arabic{enumi}.}
\tightlist
\item
  \hyperref[1-introduction-the-self-evolution-loop]{Introduction: The
  Self-Evolution Loop}
\item
  \hyperref[2-bayesian-formulation-of-gatekeeper-evolution]{Bayesian
  Formulation of Gatekeeper Evolution}
\item
  \hyperref[3-prior-distribution-over-gatekeeper-functions]{Prior
  Distribution Over Gatekeeper Functions}
\item
  \hyperref[4-likelihood-memory-bank-as-evidence]{Likelihood: Memory
  Bank as Evidence}
\item
  \hyperref[5-posterior-update-recursion]{Posterior Update Recursion}
\item
  \hyperref[6-posterior-mean-gatekeeper]{Posterior Mean Gatekeeper}
\item
  \hyperref[7-conjugate-prior-structure-gaussian-process-formulation]{Conjugate
  Prior Structure: Gaussian Process Formulation}
\item
  \hyperref[8-bayesian-martingale-property]{Bayesian Martingale
  Property}
\item
  \hyperref[9-martingale-convergence-theorem]{Martingale Convergence
  Theorem}
\item
  \hyperref[10-prior-misspecification-analysis]{Prior Misspecification
  Analysis}
\item
  \hyperref[11-kl-divergence-contraction]{KL Divergence Contraction}
\item
  \hyperref[12-bernstein-von-mises-theorem-connection]{Bernstein-von
  Mises Theorem Connection}
\item
  \hyperref[13-proof-sketches]{Proof Sketches}
\item
  \hyperref[14-summary-of-proven-vs-conjectured-claims]{Summary of
  Proven vs.~Conjectured Claims}
\item
  \hyperref[15-references]{References}
\end{enumerate}

\begin{center}\rule{0.5\linewidth}{0.5pt}\end{center}

\subsection{1. Introduction: The Self-Evolution
Loop}\label{introduction-the-self-evolution-loop}

The SCX self-evolution framework operates as a closed-loop system:

\[
\text{judge} \;\to\; \text{store} \;\to\; \text{update SCX} \;\to\; \text{re-judge} \;\to\; \text{re-update} \;\to\; \cdots
\]

At time step \(t\), the system comprises three core components:

\begin{longtable}[]{@{}
  >{\raggedright\arraybackslash}p{(\linewidth - 4\tabcolsep) * \real{0.3438}}
  >{\raggedright\arraybackslash}p{(\linewidth - 4\tabcolsep) * \real{0.2500}}
  >{\raggedright\arraybackslash}p{(\linewidth - 4\tabcolsep) * \real{0.4062}}@{}}
\toprule\noalign{}
\begin{minipage}[b]{\linewidth}\raggedright
Component
\end{minipage} & \begin{minipage}[b]{\linewidth}\raggedright
Symbol
\end{minipage} & \begin{minipage}[b]{\linewidth}\raggedright
Description
\end{minipage} \\
\midrule\noalign{}
\endhead
\bottomrule\noalign{}
\endlastfoot
Gatekeeper scoring function &
\(S_t: \mathcal{X} \times \mathcal{Y} \to [0,1]\) & Reliability score;
\(S_t(x,y)\) estimates \(P(\text{clean} \mid x, y)\) \\
Memory bank &
\(\mathcal{M}_t = \{(x_i, y_i, S_i, f_{\theta_i})\}_{i=1}^{N_t}\) &
Accumulated high-quality samples; \(N_t\) grows monotonically \\
NEP student & \(f_{\theta_t}: \mathcal{X} \to \mathcal{Y}\) &
Neural-equivalent potential trained on \(\mathcal{M}_t\) \\
\end{longtable}

The gatekeeper \(S_t\) determines which samples enter the memory bank.
The memory bank \(\mathcal{M}_t\) trains the NEP student
\(f_{\theta_t}\). The NEP student's predictions influence the next
gatekeeper update. This mutual dependence creates a coupled dynamical
system.

\textbf{Goal of this document}: Interpret the \(S_t\) evolution as a
Bayesian posterior update process, showing that self-evolution is
equivalent to sequential Bayesian inference over the space of gatekeeper
functions.

\begin{center}\rule{0.5\linewidth}{0.5pt}\end{center}

\subsection{2. Bayesian Formulation of Gatekeeper
Evolution}\label{bayesian-formulation-of-gatekeeper-evolution}

Let \(\mathcal{G} = \{S: \mathcal{X} \times \mathcal{Y} \to [0,1]\}\) be
the space of gatekeeper functions. Define a probability measure over
\(\mathcal{G}\) that represents our epistemic uncertainty about the true
data-generating process.

\textbf{Definition 4.1 (Bayesian Gatekeeper).} At time \(t\), the
gatekeeper \(S_t\) is the posterior mean of a random function
\(S \sim P_t\), where \(P_t\) is a probability distribution over
\(\mathcal{G}\):

\[
S_t(x,y) = \mathbb{E}_{S \sim P_t}[S(x,y)], \quad \forall (x,y) \in \mathcal{X} \times \mathcal{Y}
\]

The distribution \(P_t\) represents our state of knowledge at time \(t\)
about the true reliability function \(S^*\), where
\(S^*(x,y) = P(\text{clean} \mid x, y)\) is the (unknowable) oracle
reliability.

\begin{center}\rule{0.5\linewidth}{0.5pt}\end{center}

\subsection{3. Prior Distribution Over Gatekeeper
Functions}\label{prior-distribution-over-gatekeeper-functions}

\textbf{Definition 4.2 (Prior).} At initialization time \(t=0\), we
place a prior distribution \(P_0\) over \(\mathcal{G}\):

\[
P_0(S) = \text{Prior belief about the gatekeeper function}
\]

The prior may take various forms depending on available information:

\begin{enumerate}
\def\labelenumi{\arabic{enumi}.}
\tightlist
\item
  \textbf{Uninformative prior}: \(P_0(S) \propto 1\) (uniform over
  \(\mathcal{G}\)).
\item
  \textbf{SCX-initialized prior}: Using the SCX reliability scores from
  expert models:
  \[P_0(S) \propto \exp\left(-\lambda \sum_{m=1}^M \|S - SCX_m\|^2_{\mathcal{L}^2}\right)\]
  where \(SCX_m(s)\) is the state-conditioned expert reliability
  (Definition 3.1 of the existing framework), lifted to
  \(\mathcal{X} \times \mathcal{Y}\) via the state mapping.
\item
  \textbf{Gaussian process prior}: \(S \sim \mathcal{GP}(\mu_0, k_0)\)
  with mean function \(\mu_0\) and kernel \(k_0\) (see Section 7).
\end{enumerate}

The initial gatekeeper is the prior mean:
\[S_0(x,y) = \mathbb{E}_{P_0}[S(x,y)]\]

In practice, \(S_0\) may be set to the consensus score
\(C(x) = \frac{1}{M}\sum_m \mathbf{1}\{\ell(f_m(x), y) > \tau\}\) or a
constant.

\begin{center}\rule{0.5\linewidth}{0.5pt}\end{center}

\subsection{4. Likelihood: Memory Bank as
Evidence}\label{likelihood-memory-bank-as-evidence}

The memory bank
\(\mathcal{M}_t = \{(x_i, y_i, S_i, f_{\theta_i})\}_{i=1}^{N_t}\) is not
raw data but a \textbf{filtered collection}: each sample was accepted by
a previous gatekeeper \(S_i\). We must account for this selection bias.

\textbf{Definition 4.3 (Selection-Aware Likelihood).} Given a candidate
gatekeeper function \(S\), the likelihood of observing a memory bank
\(\mathcal{M}_t\) is:

\[
P(\mathcal{M}_t \mid S) = \prod_{i=1}^{N_t} P(x_i, y_i \mid S, \text{accepted at time } i)
\]

Applying Bayes' rule to the acceptance event:

\[
P(x, y \mid S, \text{accepted}) \propto P(\text{accepted} \mid x, y, S) \cdot P(x, y)
\]

The acceptance probability given \(S\) is precisely the gatekeeper
score:

\[
P(\text{accepted} \mid x, y, S) = S(x, y)
\]

\textbf{Proposition 4.1 (Likelihood Form).} Under the assumption that
the underlying data distribution \(P(x,y)\) is stationary and that
acceptance events are conditionally independent given \(S\), the
log-likelihood of the memory bank is:

\[
\log P(\mathcal{M}_t \mid S) = \sum_{i=1}^{N_t} \log S(x_i, y_i) + \sum_{i=1}^{N_t} \log P(x_i, y_i) + \text{const}
\]

where the constant accounts for the normalization of the selection
distribution.

\textbf{Remark.} The second term \(\sum \log P(x_i, y_i)\) does not
depend on \(S\) and can be absorbed into the constant. The essential
likelihood is:

\[
\log P(\mathcal{M}_t \mid S) = \sum_{i=1}^{N_t} \log S(x_i, y_i) + \text{const}(S\text{-independent})
\]

\emph{Proof sketch.} By the acceptance mechanism: each sample enters
\(\mathcal{M}_t\) only if \(S_i(x_i, y_i)\) exceeds a threshold or
otherwise passes the gate. Conditioning on the history, the probability
of observing \((x_i, y_i)\) in the bank is
\(S(x_i, y_i) \cdot P(x_i, y_i) / Z_i\), where \(Z_i\) is the acceptance
probability marginalized over the data distribution. Taking logs and
summing yields the expression. The conditional independence follows from
the i.i.d. sampling assumption (Assumption SE-A3, defined below).

\textbf{Assumption SE-A3 (Memory Bank Conditional Independence).}
Conditioned on the true gatekeeper \(S^*\), the samples in
\(\mathcal{M}_t\) are i.i.d. from the selection-biased distribution
\(P(x,y \mid S^*) \propto S^*(x,y) \cdot P(x,y)\).

This is a self-consistency condition: the data that passes the gate
should look like the data the gatekeeper considers reliable. When
iterated, this creates a virtuous (or vicious) cycle.

\begin{center}\rule{0.5\linewidth}{0.5pt}\end{center}

\subsection{5. Posterior Update
Recursion}\label{posterior-update-recursion}

\textbf{Theorem 4.1 (Bayesian Update Recursion).} Let \(P_t\) be the
posterior over gatekeeper functions at time \(t\), and let
\(\mathcal{D}_{t+1} = \{(x_j, y_j)\}_{j=1}^{B_{t+1}}\) be the new batch
of samples collected between time \(t\) and \(t+1\), each evaluated by
the current gatekeeper \(S_t\) and accepted into \(\mathcal{M}_{t+1}\)
with probability \(S_t(x_j, y_j)\). Then the posterior update is:

\[
P_{t+1}(S) \propto P_t(S) \cdot \prod_{j=1}^{B_{t+1}} S(x_j, y_j)^{a_j} \cdot (1 - S(x_j, y_j))^{1-a_j}
\]

where \(a_j \in \{0, 1\}\) is the acceptance indicator.

\textbf{Equivalent incremental form.} Writing the non-accepted data as
\(\tilde{\mathcal{D}}_{t+1}\) (samples that were evaluated but
rejected), the posterior becomes:

\[
P_{t+1}(S) \propto P_t(S) \cdot \underbrace{\left[\prod_{(x,y) \in \mathcal{D}_{t+1}} S(x,y)\right]}_{\text{accepted}} \cdot \underbrace{\left[\prod_{(x,y) \in \tilde{\mathcal{D}}_{t+1}} (1 - S(x,y))\right]}_{\text{rejected}}
\]

\emph{Interpretation.} The gatekeeper update is a Bayesian binary
classification: each sample is a Bernoulli trial with success
probability \(S(x,y)\). The accepted samples are ``successes,'' the
rejected samples are ``failures.''

\textbf{Corollary 4.1 (Posterior Recursion).} The sequence of posterior
distributions satisfies:

\[
P_{t+1}(S) \propto P_0(S) \cdot \prod_{i=1}^{t+1} \mathcal{L}_i(S \mid \mathcal{D}_i)
\]

where
\(\mathcal{L}_i(S \mid \mathcal{D}_i) = \prod_{(x,y) \in \mathcal{D}_i} S(x,y)^{a_i(x,y)} (1 - S(x,y))^{1-a_i(x,y)}\)
is the batch-\(i\) likelihood.

This is a standard sequential Bayesian update: the posterior at time
\(t\) becomes the prior for time \(t+1\).

\begin{center}\rule{0.5\linewidth}{0.5pt}\end{center}

\subsection{6. Posterior Mean
Gatekeeper}\label{posterior-mean-gatekeeper}

\textbf{Definition 4.4 (Gatekeeper as Posterior Mean).} At any time
\(t\), the deployed gatekeeper is:

\[
S_t(x,y) = \mathbb{E}_{P_t}[S(x,y)] = \int_{\mathcal{G}} S(x,y) \, dP_t(S)
\]

This is the Bayes estimator under squared error loss:

\[
S_t = \arg\min_{\hat{S}} \mathbb{E}_{P_t}[\|\hat{S} - S\|^2_{\mathcal{L}^2}]
\]

\textbf{Proposition 4.2 (Optimality of Posterior Mean).} The posterior
mean gatekeeper minimizes the expected \(\mathcal{L}^2\) prediction
error under the posterior \(P_t\):

\[
S_t = \arg\min_{\hat{S}} \mathbb{E}_{S \sim P_t} \left[ \int (S(x,y) - \hat{S}(x,y))^2 \, dP(x,y) \right]
\]

\emph{Proof.} Standard result: the posterior mean minimizes Bayes risk
under squared error loss. See {[}Berger, 1985, Section 4.4{]}.

\textbf{Connection to SCX consensus.} When the prior is flat and only
the initial expert-based SCX scores are used, the posterior mean reduces
to:

\[
S_0(x,y) \approx \frac{1}{M}\sum_{m=1}^M \mathbf{1}\{\ell(f_m(x), y) < \tau\}
\]

which recovers the consensus score \(1 - C(x)\) from Theorem 1. This
bridges the Bayesian formulation to the existing SCX theory.

\begin{center}\rule{0.5\linewidth}{0.5pt}\end{center}

\subsection{7. Conjugate Prior Structure: Gaussian Process
Formulation}\label{conjugate-prior-structure-gaussian-process-formulation}

\textbf{Definition 4.5 (Gaussian Process Gatekeeper).} Let the
gatekeeper function \(S: \mathcal{X} \times \mathcal{Y} \to \mathbb{R}\)
be modeled as a Gaussian process (GP) prior:

\[
S \sim \mathcal{GP}(\mu_0, k_0)
\]

where \(\mu_0: \mathcal{Z} \to \mathbb{R}\) is the mean function
(\(\mathcal{Z} = \mathcal{X} \times \mathcal{Y}\)) and
\(k_0: \mathcal{Z} \times \mathcal{Z} \to \mathbb{R}\) is the covariance
kernel.

\textbf{Assumption SE-A4 (Logit-Normal Likelihood).} The acceptance
probability is modeled via a logistic link:

\[
P(a = 1 \mid S, x, y) = \sigma(S(x,y)) = \frac{1}{1 + e^{-S(x,y)}}
\]

where \(\sigma\) is the logistic sigmoid.

\textbf{Theorem 4.2 (GP Posterior Update).} Under the GP prior and
logistic likelihood, the posterior at time \(t\) is:

\[
P_t(S \mid \mathcal{M}_t) \propto \mathcal{GP}(\mu_0, k_0) \cdot \prod_{i=1}^{N_t} \sigma(S(x_i, y_i))^{a_i} (1 - \sigma(S(x_i, y_i)))^{1-a_i}
\]

While the exact posterior remains non-Gaussian (due to the logistic
likelihood), the \textbf{Laplace approximation} yields a Gaussian
posterior:

\[
P_t(S) \approx \mathcal{GP}(\mu_t, k_t)
\]

where:

\begin{itemize}
\tightlist
\item
  \(\mu_t(z) = \mu_0(z) + k_0(z, \mathbf{Z}_t)^\top (k_0(\mathbf{Z}_t, \mathbf{Z}_t) + W_t^{-1})^{-1} (\mathbf{a}_t - \sigma(\mu_0(\mathbf{Z}_t)))\)
\item
  \(k_t(z, z') = k_0(z, z') - k_0(z, \mathbf{Z}_t)^\top (k_0(\mathbf{Z}_t, \mathbf{Z}_t) + W_t^{-1})^{-1} k_0(\mathbf{Z}_t, z')\)
\item
  \(W_t\) is a diagonal matrix with
  \(W_{t,ii} = \sigma(\mu_t(z_i))(1 - \sigma(\mu_t(z_i)))\)
\item
  \(\mathbf{Z}_t\) is the concatenated set of all evaluated points
  \(((x_i, y_i))_{i=1}^{N_t}\)
\item
  \(\mathbf{a}_t\) is the vector of acceptance indicators
\end{itemize}

\textbf{Closed-form posterior mean.} Under the Laplace approximation:

\[
S_t(z) = \mu_t(z) = \mathbb{E}_{P_t}[S(z)] \approx k_0(z, \mathbf{Z}_t)^\top (k_0(\mathbf{Z}_t, \mathbf{Z}_t) + W_t^{-1})^{-1} (\mathbf{a}_t - \sigma(\mu_0(\mathbf{Z}_t)) + W_t \mu_0(\mathbf{Z}_t))
\]

\textbf{Conjugate special case: Gaussian likelihood.} If we use a
Gaussian observation model (acceptance score is continuous in \([0,1]\)
rather than binary), the prior and posterior are both Gaussian:

\[
S_t(z) = k_0(z, \mathbf{Z}_t)^\top (k_0(\mathbf{Z}_t, \mathbf{Z}_t) + \sigma^2_\varepsilon I)^{-1} \mathbf{a}_t
\]

This is the standard GP regression form, with closed-form sequential
update.

\textbf{Proposition 4.3 (Conjugate Sequential GP Update).} Under the
Gaussian likelihood model, the posterior can be updated sequentially
without revisiting old data:

\[
\mu_{t+1}(z) = \mu_t(z) + k_t(z, z_{t+1}) \cdot \frac{a_{t+1} - \mu_t(z_{t+1})}{k_t(z_{t+1}, z_{t+1}) + \sigma^2_\varepsilon}
\]

\[
k_{t+1}(z, z') = k_t(z, z') - \frac{k_t(z, z_{t+1}) k_t(z_{t+1}, z')}{k_t(z_{t+1}, z_{t+1}) + \sigma^2_\varepsilon}
\]

\emph{Proof.} Standard sequential GP update formulas. See {[}Rasmussen
\& Williams, 2006, Algorithm 2.1{]}. The conjugate property holds
because the Gaussian-Gaussian pair forms a conditionally conjugate
family.

\begin{center}\rule{0.5\linewidth}{0.5pt}\end{center}

\subsection{8. Bayesian Martingale
Property}\label{bayesian-martingale-property}

\textbf{Theorem 4.3 (Posterior Mean Martingale).} Let
\(S_t(x,y) = \mathbb{E}_{P_t}[S(x,y)]\) be the posterior mean
gatekeeper, where \(P_t\) is updated via Bayes' rule as new batches
arrive. Then the stochastic process \(\{S_t\}_{t \geq 0}\) is a
martingale with respect to the filtration
\(\mathcal{F}_t = \sigma(\mathcal{M}_1, \dots, \mathcal{M}_t)\):

\[
\mathbb{E}[S_{t+1}(x,y) \mid \mathcal{F}_t] = S_t(x,y), \quad \forall (x,y) \in \mathcal{X} \times \mathcal{Y}
\]

\emph{Proof.} Apply the tower property of conditional expectation:

\[
\begin{aligned}
\mathbb{E}[S_{t+1}(x,y) \mid \mathcal{F}_t] &= \mathbb{E}[\mathbb{E}_{P_{t+1}}[S(x,y)] \mid \mathcal{F}_t] \\
&= \mathbb{E}[\mathbb{E}[S(x,y) \mid \mathcal{M}_{1:t+1}] \mid \mathcal{F}_t] \\
&= \mathbb{E}[S(x,y) \mid \mathcal{M}_{1:t}] \quad \text{(tower property)} \\
&= \mathbb{E}_{P_t}[S(x,y)] = S_t(x,y)
\end{aligned}
\]

The key step is that \(P_{t+1}\) is the posterior after conditioning on
\(\mathcal{M}_{t+1}\), and the tower property of conditional expectation
gives
\(\mathbb{E}[\mathbb{E}[S \mid \mathcal{F}_{t+1}] \mid \mathcal{F}_t] = \mathbb{E}[S \mid \mathcal{F}_t]\).

\textbf{Corollary 4.2 (Pointwise Martingale).} For any fixed \((x,y)\),
the sequence \(\{S_t(x,y)\}_{t \geq 0}\) is a martingale bounded in
\([0,1]\).

\textbf{Corollary 4.3 (Vector Martingale).} If
\(\mathcal{Z}_0 \subset \mathcal{X} \times \mathcal{Y}\) is a finite set
of evaluation points, then \(\{S_t(z)\}_{t \geq 0}\) for
\(z \in \mathcal{Z}_0\) is a vector-valued martingale.

\begin{center}\rule{0.5\linewidth}{0.5pt}\end{center}

\subsection{9. Martingale Convergence
Theorem}\label{martingale-convergence-theorem}

\textbf{Theorem 4.4 (Martingale Convergence of Gatekeeper).} Under the
Bayesian update framework, the sequence of gatekeeper functions
\(\{S_t\}_{t \geq 0}\) converges almost surely to a limit \(S_\infty\):

\[
S_t(x,y) \xrightarrow{a.s.} S_\infty(x,y), \quad \forall (x,y) \in \mathcal{X} \times \mathcal{Y}
\]

as \(t \to \infty\), where \(S_\infty\) is a random function measurable
with respect to
\(\mathcal{F}_\infty = \sigma(\bigcup_{t \geq 0} \mathcal{F}_t)\).

\emph{Proof.} The martingale convergence theorem (Doob, 1953) states
that any martingale bounded in \(\mathcal{L}^1\) converges almost
surely. Since \(S_t(x,y) \in [0,1]\) for all \(t\), it is uniformly
bounded and thus \(\mathcal{L}^1\)-bounded. Therefore
\(S_t(x,y) \to S_\infty(x,y)\) almost surely. By the separability of
\(\mathcal{X} \times \mathcal{Y}\) (under mild topological conditions),
convergence holds simultaneously for all \((x,y)\) in a dense subset,
and by continuity of the mapping \(z \mapsto S_t(z)\) (if the kernel
\(k_0\) is continuous), convergence is uniform on compact sets. See Doob
(1953), Chapter VII.

\textbf{Proposition 4.4 (Structure of the Limit).} The limit
\(S_\infty\) satisfies the self-consistency equation:

\[
S_\infty(x,y) = \mathbb{E}[S(x,y) \mid \mathcal{F}_\infty] = \int S(x,y) \, dP_\infty(S)
\]

where \(P_\infty\) is the limiting posterior distribution. Moreover,
\(S_\infty\) is the Bayes estimator under the limiting posterior:

\[
S_\infty = \arg\min_{\hat{S}} \mathbb{E}_{P_\infty}[\|\hat{S} - S\|^2_{\mathcal{L}^2}]
\]

\textbf{Proposition 4.5 (Convergence Rate).} Under the GP prior with
Gaussian likelihood, the convergence rate is:

\[
\|S_t - S_\infty\|_{\mathcal{L}^2} = O_p(t^{-1/2})
\]

if the covariance kernel \(k_0\) has finite trace and the observation
noise \(\sigma^2_\varepsilon\) is bounded away from zero.

\emph{Proof sketch.} For GP regression with i.i.d. Gaussian noise, the
posterior variance at any point decays as \(O(1/N_t)\). Since
\(N_t \to \infty\) as \(t \to \infty\) (memory bank grows), the
posterior mean converges to the true function at rate
\(O_p(1/\sqrt{N_t})\). If \(N_t \propto t\) (constant batch size), the
rate is \(O_p(t^{-1/2})\). A more precise rate depends on the kernel's
eigenvalue decay; for a Matérn kernel with smoothness \(\nu\), the rate
is \(O_p(t^{-\nu/(2\nu+d)})\).

\begin{center}\rule{0.5\linewidth}{0.5pt}\end{center}

\subsection{10. Prior Misspecification
Analysis}\label{prior-misspecification-analysis}

\textbf{Definition 4.6 (Misspecified Prior).} The prior \(P_0\) is
\textbf{misspecified} if the true oracle gatekeeper \(S^*\) lies outside
the support of \(P_0\):

\[
S^* \notin \text{supp}(P_0)
\]

\textbf{Theorem 4.5 (Behavior Under Prior Misspecification).} Let
\(S^*\) be the true oracle gatekeeper, and let \(P_0\) be a prior whose
support is a set \(\mathcal{G}_0 \subset \mathcal{G}\). If
\(S^* \notin \mathcal{G}_0\), then:

\textbf{(a) Posterior inconsistency.} The posterior distribution \(P_t\)
does \textbf{not} concentrate on \(S^*\):

\[
P_t(\{S: \|S - S^*\| > \varepsilon\}) \not\to 0
\]

for any \(\varepsilon < \inf_{S \in \mathcal{G}_0} \|S - S^*\|\).

\textbf{(b) Convergence to KL projection.} Instead, \(P_t\) concentrates
on the set of functions in \(\mathcal{G}_0\) that minimize KL divergence
to the true data-generating process:

\[
P_t \to \delta_{S^\dagger}, \quad S^\dagger = \arg\min_{S \in \mathcal{G}_0} D_{\text{KL}}(P^* \parallel P_S)
\]

where \(P^*\) is the true data distribution and \(P_S\) is the
distribution induced by \(S \in \mathcal{G}_0\).

\textbf{(c) Gatekeeper limit.} The limiting gatekeeper \(S_\infty\) is
the posterior mean under \(P_\infty\), which converges to:

\[
S_\infty = \mathbb{E}_{P_\infty}[S] = S^\dagger
\]

where \(S^\dagger\) is the KL projection of \(S^*\) onto
\(\mathcal{G}_0\).

\emph{Proof sketch.} Part (a) follows from the definition of support: if
\(S^* \notin \mathcal{G}_0\), then for any \(\varepsilon\) smaller than
the distance from \(S^*\) to \(\mathcal{G}_0\), the posterior
probability of the \(\varepsilon\)-ball around \(S^*\) is zero at all
\(t\), so it cannot converge to 1. Part (b) follows from the
Bernstein-von Mises theorem for misspecified models (Kleijn \& van der
Vaart, 2012): the posterior concentrates on the set of distributions
that minimize KL divergence to the truth. Part (c) follows from
continuous mapping: if \(P_t \to \delta_{S^\dagger}\), then
\(\mathbb{E}_{P_t}[S] \to S^\dagger\).

\textbf{Corollary 4.4 (Practical Implications of Misspecification).}

\begin{enumerate}
\def\labelenumi{\arabic{enumi}.}
\item
  \textbf{SCX-Expert prior}: If \(S_0\) is initialized from expert
  consensus scores and the experts are systematically biased, \(S^*\)
  may be outside the manifold of functions realizable by the
  expert-based prior. The gatekeeper converges to the best approximation
  within the prior family, not the truth.
\item
  \textbf{GP with wrong kernel}: If \(S^*\) has a different smoothness
  or length scale than assumed by \(k_0\), the posterior mean converges
  to the best GP approximation of \(S^*\), which may have non-vanishing
  bias.
\item
  \textbf{Detection of misspecification}: The martingale property
  provides a diagnostic: if \(S_t\) shows systematic drift beyond what
  the posterior variance predicts, it suggests prior misspecification:
  \[
  \mathbb{E}[(S_{t+1} - S_t)^2] > \mathbb{E}_t[\text{Var}_{P_t}[S]]
  \]
\end{enumerate}

\begin{center}\rule{0.5\linewidth}{0.5pt}\end{center}

\subsection{11. KL Divergence
Contraction}\label{kl-divergence-contraction}

\textbf{Definition 4.7 (KL Divergence Between Data Distributions).} Let
\(P^*\) be the true data-generating distribution induced by the oracle
gatekeeper \(S^*\), and let \(P_t\) be the distribution induced by the
posterior predictive:

\[
P_t(x,y) = \int S(x,y) \cdot P(x,y) \, dP_t(S)
\]

The KL divergence from the truth to the model at time \(t\) is:

\[
D_{\text{KL}}(P^* \parallel P_t) = \mathbb{E}_{(x,y) \sim P^*} \left[ \log \frac{P^*(x,y)}{P_t(x,y)} \right]
\]

\textbf{Theorem 4.6 (KL Contraction).} Under the Bayesian update
framework with correctly specified prior (\(S^* \in \text{supp}(P_0)\))
and regularity conditions (identifiability, integrability, and Doob's
consistency conditions):

\[
D_{\text{KL}}(P^* \parallel P_t) \xrightarrow{a.s.} 0 \quad \text{as } t \to \infty
\]

Moreover, the contraction rate satisfies:

\[
\mathbb{E}[D_{\text{KL}}(P^* \parallel P_t)] \leq \frac{D_{\text{KL}}(P^* \parallel P_0)}{t} \quad \text{(under i.i.d. sampling)}
\]

\emph{Proof sketch.} This follows from Doob's consistency theorem for
Bayesian models (Doob, 1949; Ghosh \& Ramamoorthi, 2003). The key
conditions are: 1. The model is identifiable:
\(S \neq S' \implies P_S \neq P_{S'}\). 2. The prior puts positive mass
on all neighborhoods of \(S^*\). 3. The data are i.i.d. conditional on
\(S\).

Under these conditions, the posterior is consistent at \(S^*\) for
almost every \(S^*\) in the support of the prior, and KL divergence
contracts to zero. The rate inequality follows from the convexity of KL
divergence and the martingale property of the log-likelihood ratio.

\textbf{Proposition 4.6 (KL Decomposition).} The KL divergence
decomposes into prior mismatch and learning components:

\[
D_{\text{KL}}(P^* \parallel P_t) = \underbrace{D_{\text{KL}}(P^* \parallel P_0)}_{\text{initial gap}} - \underbrace{\sum_{i=1}^t \mathbb{E}_{P^*}[\log \ell_i(S) - \log \mathbb{E}_{P_{i-1}}[\ell_i(S)]]}_{\text{cumulative learning}}
\]

where \(\ell_i(S) = P(\mathcal{D}_i \mid S)\) is the likelihood of batch
\(i\).

\emph{Proof.} Write the posterior
\(P_t \propto P_0 \prod_{i=1}^t \ell_i(S)\). Then
\(\log dP_t/dP_0 = \sum_{i=1}^t \log \ell_i(S) - \log Z_t\) where
\(Z_t\) is the marginal likelihood. Taking expectations and telescoping
yields the decomposition.

\begin{center}\rule{0.5\linewidth}{0.5pt}\end{center}

\subsection{12. Bernstein-von Mises Theorem
Connection}\label{bernstein-von-mises-theorem-connection}

\textbf{Theorem 4.7 (Bernstein-von Mises for Gatekeeper).} Under
regularity conditions (identifiability, smoothness, and correct
specification), the posterior distribution \(P_t\) converges to a
Gaussian process centered at the truth \(S^*\) with covariance
\(N_t^{-1} \mathcal{I}(S^*)^{-1}\):

\[
\sqrt{N_t} (P_t - \delta_{S^*}) \xrightarrow{d} \mathcal{GP}(0, \mathcal{I}(S^*)^{-1})
\]

in the sense of convergence of posterior distributions. Equivalently,
for any bounded continuous functional \(f: \mathcal{G} \to \mathbb{R}\):

\[
\sqrt{N_t} \left( \mathbb{E}_{P_t}[f(S)] - f(S^*) \right) \xrightarrow{d} \mathcal{N}(0, \nabla f(S^*)^\top \mathcal{I}(S^*)^{-1} \nabla f(S^*))
\]

\emph{Conditions.} The theorem requires:

\begin{enumerate}
\def\labelenumi{\arabic{enumi}.}
\tightlist
\item
  \textbf{Identifiability}: \(S \neq S^* \implies P_S \neq P_{S^*}\).
\item
  \textbf{Smoothness}: The log-likelihood
  \(\log P(\mathcal{M}_t \mid S)\) is twice continuously differentiable
  in a neighborhood of \(S^*\).
\item
  \textbf{Positivity}: The Fisher information
  \(\mathcal{I}(S) = -\mathbb{E}_{P^*}[\nabla^2 \log P(\mathcal{M}_1 \mid S)]\)
  is positive definite and continuous at \(S^*\).
\item
  \textbf{Prior regularity}: The prior \(P_0\) has a density (with
  respect to an appropriate base measure) that is continuous and
  positive at \(S^*\).
\item
  \textbf{Memory bank size}: \(N_t \to \infty\).
\end{enumerate}

\emph{Proof sketch.} The Bernstein-von Mises theorem (van der Vaart,
1998, Chapter 10) states that under regularity conditions, the posterior
distribution is asymptotically normal. The infinite-dimensional
extension (Gaussian process version) follows from the functional central
limit theorem for posterior distributions (Castillo \& Nickl, 2013). The
key steps are: (1) expand the log-likelihood around \(S^*\), (2) show
that the prior becomes negligible as \(N_t \to \infty\), (3) identify
the limiting normal distribution via the Fisher information.

\textbf{Corollary 4.5 (Asymptotic Normality of Gatekeeper).} For any
fixed \((x,y)\), as \(t \to \infty\):

\[
\sqrt{N_t} (S_t(x,y) - S^*(x,y)) \xrightarrow{d} \mathcal{N}(0, \mathcal{I}(S^*)^{-1}(x,y))
\]

where \(\mathcal{I}(S^*)^{-1}(x,y)\) is the pointwise posterior
variance.

\textbf{Corollary 4.6 (Asymptotic Credible Intervals).} Under the
Bernstein-von Mises theorem, the \(1-\alpha\) credible interval for
\(S^*(x,y)\) is asymptotically:

\[
S_t(x,y) \pm z_{\alpha/2} \cdot \sqrt{\frac{\hat{\sigma}^2_t(x,y)}{N_t}}
\]

where \(\hat{\sigma}^2_t(x,y)\) is the estimated posterior variance and
\(z_{\alpha/2}\) is the standard normal quantile.

\textbf{Important caveat.} The Bernstein-von Mises theorem requires
correct specification (\(S^*\) in the support of \(P_0\)). Under
misspecification (Section 10), the posterior converges to a Gaussian
centered at \(S^\dagger\) (the KL projection), not \(S^*\). This is the
\textbf{misspecified Bernstein-von Mises} phenomenon (Kleijn \& van der
Vaart, 2012).

\begin{center}\rule{0.5\linewidth}{0.5pt}\end{center}

\subsection{13. Proof Sketches}\label{proof-sketches}

\subsubsection{Proof Sketch for Theorem 4.3 (Martingale
Property)}\label{proof-sketch-for-theorem-4.3-martingale-property}

\begin{enumerate}
\def\labelenumi{\arabic{enumi}.}
\item
  \textbf{Setup.} Let \((\Omega, \mathcal{F}, \mathbb{P})\) be the
  underlying probability space. Define
  \(\mathcal{F}_t = \sigma(\mathcal{M}_1, \dots, \mathcal{M}_t)\) as the
  filtration generated by the memory bank up to time \(t\).
\item
  \textbf{Conditional expectation.} By definition,
  \(S_t(x,y) = \mathbb{E}[S(x,y) \mid \mathcal{F}_t]\) where
  \(S \sim P_0\) is the random function drawn from the prior, and the
  expectation is over both \(S\) and the data.
\item
  \textbf{Tower property.} For any \(t\): \[
  \mathbb{E}[S_{t+1}(x,y) \mid \mathcal{F}_t] = \mathbb{E}[\mathbb{E}[S(x,y) \mid \mathcal{F}_{t+1}] \mid \mathcal{F}_t] = \mathbb{E}[S(x,y) \mid \mathcal{F}_t] = S_t(x,y)
  \] The middle equality is the tower property of conditional
  expectation.
\item
  \textbf{Integrability.} Since \(S_t(x,y) \in [0,1]\), we have
  \(\mathbb{E}[|S_t(x,y)|] \leq 1 < \infty\), so the martingale property
  holds.
\end{enumerate}

\subsubsection{Proof Sketch for Theorem 4.4 (Martingale
Convergence)}\label{proof-sketch-for-theorem-4.4-martingale-convergence}

\begin{enumerate}
\def\labelenumi{\arabic{enumi}.}
\item
  \textbf{Bounded martingale.} \(\{S_t(x,y)\}_{t \geq 0}\) is a
  martingale bounded in \([0,1]\).
\item
  \textbf{Doob's convergence theorem.} Any martingale that is bounded in
  \(\mathcal{L}^1\) converges almost surely. Since \(S_t\) is uniformly
  bounded, \(\sup_t \mathbb{E}[|S_t|] \leq 1 < \infty\), so
  \(\mathcal{L}^1\)-boundedness holds.
\item
  \textbf{Almost sure convergence.} Hence \(S_t(x,y) \to S_\infty(x,y)\)
  almost surely for each fixed \((x,y)\).
\item
  \textbf{Simultaneous convergence.} For a countable dense subset
  \(\mathcal{Z}_0 \subset \mathcal{X} \times \mathcal{Y}\), convergence
  holds simultaneously for all \(z \in \mathcal{Z}_0\) with probability
  1. If the sample paths \(t \mapsto S_t(z)\) are continuous in \(z\)
  (guaranteed by the GP kernel), then convergence extends uniformly on
  compact sets.
\end{enumerate}

\subsubsection{Proof Sketch for Proposition 4.1 (Likelihood
Form)}\label{proof-sketch-for-proposition-4.1-likelihood-form}

\begin{enumerate}
\def\labelenumi{\arabic{enumi}.}
\item
  \textbf{Acceptance mechanism.} Sample \((x,y)\) is accepted into
  \(\mathcal{M}_t\) with probability \(S(x,y)\), independently of other
  samples.
\item
  \textbf{Selection distribution.} The distribution of accepted samples
  is: \[
  P(x,y \mid \text{accepted}, S) = \frac{S(x,y) \cdot P(x,y)}{\int S(x',y') \cdot P(x',y') \, dx' dy'}
  \]
\item
  \textbf{Joint likelihood.} For the entire memory bank
  \(\mathcal{M}_t = \{(x_i, y_i)\}_{i=1}^{N_t}\): \[
  P(\mathcal{M}_t \mid S) = \prod_{i=1}^{N_t} \frac{S(x_i, y_i) \cdot P(x_i, y_i)}{Z(S)}
  \] where \(Z(S) = \mathbb{E}_{P}[S(X,Y)]\) is the marginal acceptance
  probability.
\item
  \textbf{Log-likelihood.} Taking logs: \[
  \log P(\mathcal{M}_t \mid S) = \sum_{i=1}^{N_t} \log S(x_i, y_i) + \sum_{i=1}^{N_t} \log P(x_i, y_i) - N_t \log Z(S)
  \]
\item
  \textbf{Constant elimination.} The term \(\sum \log P(x_i, y_i)\) does
  not depend on \(S\) and can be treated as constant. The term
  \(N_t \log Z(S)\) is constant across \(S\) if \(Z(S)\) is
  approximately constant (which holds when the prior is broad and the
  acceptance threshold is near-independent of \(S\)). In the full
  Bayesian treatment, this term contributes to the evidence
  normalization.
\end{enumerate}

\subsubsection{Proof Sketch for Theorem 4.6 (KL
Contraction)}\label{proof-sketch-for-theorem-4.6-kl-contraction}

\begin{enumerate}
\def\labelenumi{\arabic{enumi}.}
\item
  \textbf{Prior support condition.} By assumption,
  \(S^* \in \text{supp}(P_0)\), meaning
  \(P_0(\{S: \|S - S^*\| < \varepsilon\}) > 0\) for all
  \(\varepsilon > 0\).
\item
  \textbf{Doob's consistency.} For almost every \(S\) in the support of
  \(P_0\), the posterior is consistent (Doob, 1949). Since \(S^*\) is in
  the support (by correct specification), consistency holds.
\item
  \textbf{KL convergence.} Posteriour consistency implies that \(P_t\)
  concentrates on \(S^*\), which means \(P_t(x,y) \to P^*(x,y)\) in
  total variation, and hence in KL divergence (by Pinsker's inequality
  and the fact that convergence in TV implies convergence in KL when the
  limit density is bounded away from zero).
\item
  \textbf{Rate.} Under i.i.d. sampling and the LAN (local asymptotic
  normality) condition, the rate of KL contraction is \(O(1/t)\).
  Specifically: \[
  \mathbb{E}[D_{\text{KL}}(P^* \parallel P_t)] \leq \frac{D_{\text{KL}}(P^* \parallel P_0)}{t}
  \] follows from the convexity of KL divergence and the fact that the
  average of \(t\) conditionally independent updates gives a telescoping
  bound.
\end{enumerate}

\begin{center}\rule{0.5\linewidth}{0.5pt}\end{center}

\subsection{14. Summary of Proven vs.~Conjectured
Claims}\label{summary-of-proven-vs.-conjectured-claims}

\begin{longtable}[]{@{}
  >{\raggedright\arraybackslash}p{(\linewidth - 4\tabcolsep) * \real{0.3182}}
  >{\raggedright\arraybackslash}p{(\linewidth - 4\tabcolsep) * \real{0.3636}}
  >{\raggedright\arraybackslash}p{(\linewidth - 4\tabcolsep) * \real{0.3182}}@{}}
\toprule\noalign{}
\begin{minipage}[b]{\linewidth}\raggedright
Claim
\end{minipage} & \begin{minipage}[b]{\linewidth}\raggedright
Status
\end{minipage} & \begin{minipage}[b]{\linewidth}\raggedright
Notes
\end{minipage} \\
\midrule\noalign{}
\endhead
\bottomrule\noalign{}
\endlastfoot
Theorem 4.1: Bayesian update recursion & \textbf{Proven} & Direct
application of Bayes' rule \\
Proposition 4.1: Likelihood form & \textbf{Proven} & Derived from
acceptance mechanism \\
Proposition 4.2: Posterior mean optimality & \textbf{Proven} & Standard
Bayes risk minimization \\
Theorem 4.3: Posterior mean martingale & \textbf{Proven} & Tower
property + conditional expectation \\
Corollary 4.2: Pointwise martingale & \textbf{Proven} & Direct
consequence of Theorem 4.3 \\
Theorem 4.4: Martingale convergence & \textbf{Proven} & Doob's
martingale convergence theorem \\
Proposition 4.3: Conjugate GP update & \textbf{Proven} & Standard GP
regression formulas \\
Theorem 4.5: Prior misspecification & \textbf{Proven} & KL projection +
posterior concentration theory \\
Theorem 4.6: KL contraction & \textbf{Proven under regularity
conditions} & Doob's consistency + LAN conditions \\
Proposition 4.4: Structure of the limit & \textbf{Proven} & Limit of
Bayes estimator under squared error \\
Proposition 4.5: Convergence rate & \textbf{Conjectured} & Rate
\(O_p(t^{-1/2})\) holds for GP regression with finite-trace kernel;
general case depends on kernel eigen-decay \\
Theorem 4.7: Bernstein-von Mises & \textbf{Proven under regularity
conditions} & Standard BvM + infinite-dimensional extension (Castillo \&
Nickl, 2013) \\
Proposition 4.6: KL decomposition & \textbf{Proven} & Algebraic
manipulation of log posterior \\
Corollary 4.4: Misspecification diagnostics & \textbf{Heuristic} & The
drift diagnostic is empirically motivated \\
\end{longtable}

\textbf{Key uncertainty}: The convergence rate (Proposition 4.5) depends
on the specific kernel choice and the eigenstructure of the covariance
operator. For Matérn kernels, the rate is known; for deep kernel GPs
(which may be more realistic for SCX), rates are an active research
area.

\begin{center}\rule{0.5\linewidth}{0.5pt}\end{center}

\subsection{15. References}\label{references}

\begin{enumerate}
\def\labelenumi{\arabic{enumi}.}
\tightlist
\item
  Doob, J. L. (1949). Application of the theory of martingales to the
  theory of Markov chains. \emph{Colloques Internationaux du CNRS}, 13,
  23-27.
\item
  Doob, J. L. (1953). \emph{Stochastic Processes}. Wiley.
\item
  Berger, J. O. (1985). \emph{Statistical Decision Theory and Bayesian
  Analysis} (2nd ed.). Springer.
\item
  Ghosh, J. K., \& Ramamoorthi, R. V. (2003). \emph{Bayesian
  Nonparametrics}. Springer.
\item
  van der Vaart, A. W. (1998). \emph{Asymptotic Statistics}. Cambridge
  University Press.
\item
  Kleijn, B. J. K., \& van der Vaart, A. W. (2012). The Bernstein-von
  Mises theorem under misspecification. \emph{Electronic Journal of
  Statistics}, 6, 354-381.
\item
  Castillo, I., \& Nickl, R. (2013). Nonparametric Bernstein-von Mises
  theorems in Gaussian white noise. \emph{Annals of Statistics}, 41(4),
  1999-2028.
\item
  Rasmussen, C. E., \& Williams, C. K. I. (2006). \emph{Gaussian
  Processes for Machine Learning}. MIT Press.
\item
  Kleijn, B. J. K. (2020). \emph{Bayesian Theory with Applications}.
  Cambridge University Press.
\item
  Dawid, A. P. (1984). Statistical theory: The prequential approach.
  \emph{JRSS Series A}, 147(2), 278-292.
\item
  Williams, C. K. I., \& Rasmussen, C. E. (1996). Gaussian processes for
  regression. \emph{NIPS 1995}.
\item
  O'Hagan, A. (1978). Curve fitting and optimal design for prediction.
  \emph{JRSS Series B}, 40(1), 1-42.
\end{enumerate}

\begin{center}\rule{0.5\linewidth}{0.5pt}\end{center}

\emph{End of 04\_bayesian\_update.md -- Bayesian update interpretation
of SCX self-evolution.}

\end{document}
